% !TEX program = xelatex
% ============================================================
%  قائمة مالية — Financial Statement Template
%  قائمة الدخل وقائمة المركز المالي باللغة العربية
%  Compile with: xelatex main.tex (run twice)
% ============================================================
%  Features:
%    - Header with company name, period, and statement title
%    - Income statement table (Revenue, COGS, Gross profit, Opex, Net income)
%    - Balance sheet summary (Assets, Liabilities, Equity)
%    - Two-column layout (Item | Amount) with booktabs
% ============================================================

\documentclass[11pt,a4paper]{article}

\usepackage[a4paper,margin=2cm,top=2cm,bottom=2cm,headheight=16pt]{geometry}
\usepackage{graphicx}
\usepackage{array}
\usepackage{booktabs}
\usepackage{tabularx}
\usepackage{multirow}
\usepackage{enumitem}
\usepackage{float}
\usepackage{titlesec}
\usepackage{fancyhdr}
\usepackage{setspace}
\usepackage[table]{xcolor}

\usepackage{bidi}
\usepackage[no-math]{fontspec}
\usepackage[quiet]{polyglossia}

\setmainfont[Scale=1]{NotoKufiArabic-Regular.ttf}
\newfontfamily\arabicfont[Script=Arabic,Scale=0.95]{NotoKufiArabic-Regular.ttf}
\newfontfamily\englishfont[Scale=0.95]{TeX Gyre Termes}

\setdefaultlanguage[calendar=gregorian,hijricorrection=1]{arabic}
\setotherlanguage[variant=british]{english}

\usepackage[breaklinks,hidelinks]{hyperref}
\hypersetup{pdftitle={القائمة المالية},pdfauthor={اسم الشركة}}

% Currency command — TODO: change currency code as needed (ر.س = SAR)
\newcommand{\currency}[1]{#1\,\textarabic{ر.س}}

\titleformat{\section}{\large\bfseries}{\thesection}{0.7em}{}
\titlespacing*{\section}{0pt}{1.2ex plus 0.3ex}{0.8ex plus 0.2ex}

\pagestyle{fancy}
\fancyhf{}
% TODO: set period placeholder
\fancyhead[R]{\small\itshape القائمة المالية (\LR{Financial Statement})}
\fancyhead[L]{\small اسم الشركة}
\fancyfoot[C]{\small\thepage}
\renewcommand{\headrulewidth}{0.3pt}

\newcolumntype{L}[1]{>{\raggedright\arraybackslash}p{#1}}
\newcolumntype{C}[1]{>{\centering\arraybackslash}p{#1}}
\newcolumntype{R}[1]{>{\raggedleft\arraybackslash}p{#1}}

\definecolor{accentColor}{HTML}{1A3C6B}
\definecolor{revenueColor}{HTML}{EAF7EE}
\definecolor{expenseColor}{HTML}{FBEAEA}
\definecolor{totalColor}{HTML}{EAF0F6}

\setstretch{1.2}

\begin{document}

% ============================================================
% HEADER
% ============================================================
\begin{center}
{\LARGE\bfseries\color{accentColor} اسم الشركة}\\[0.3em]
{\large القائمة المالية (\LR{Financial Statement})}\\[0.3em]
{\normalsize للفترة من \rule{2.5cm}{0.4pt} إلى \rule{2.5cm}{0.4pt}} % TODO: set period
\\[1em]
\end{center}

\vspace{0.5em}

% ============================================================
% 1. INCOME STATEMENT
% ============================================================
\section*{قائمة الدخل (\LR{Income Statement})}
% TODO: replace sample amounts with actual figures

\begin{table}[H]
\centering
\renewcommand{\arraystretch}{1.4}
\begin{tabularx}{\textwidth}{L{10cm} R{5cm}}
\toprule
\textbf{البند (\LR{Item})} & \textbf{المبلغ (\LR{Amount})} \\
\midrule
\rowcolor{revenueColor}
\multicolumn{2}{L{15cm}}{\textbf{الإيرادات (\LR{Revenue})}} \\
\midrule
إيرادات المبيعات (\LR{Sales Revenue}) & \currency{850,000.00} \\
إيرادات الخدمات (\LR{Service Revenue}) & \currency{320,000.00} \\
إيرادات أخرى (\LR{Other Income}) & \currency{30,000.00} \\
\midrule
\textbf{إجمالي الإيرادات (\LR{Total Revenue})} & \textbf{\currency{1,200,000.00}} \\
\midrule
\rowcolor{expenseColor}
\multicolumn{2}{L{15cm}}{\textbf{تكلفة البضاعة المباعة (\LR{Cost of Goods Sold})}} \\
\midrule
المواد الخام (\LR{Raw Materials}) & \currency{280,000.00} \\
الأجور المباشرة (\LR{Direct Labor}) & \currency{150,000.00} \\
مصاريف إنتاج (\LR{Manufacturing Overhead}) & \currency{90,000.00} \\
\midrule
\textbf{إجمالي تكلفة البضاعة المباعة (\LR{Total COGS})} & \textbf{\currency{520,000.00}} \\
\midrule
\rowcolor{totalColor}
\textbf{\large مجمل الربح (\LR{Gross Profit})} & \textbf{\large\color{accentColor}\currency{680,000.00}} \\
\midrule
\rowcolor{expenseColor}
\multicolumn{2}{L{15cm}}{\textbf{المصاريف التشغيلية (\LR{Operating Expenses})}} \\
\midrule
الرواتب والإيجارات (\LR{Salaries \& Rent}) & \currency{180,000.00} \\
التسويق والإعلان (\LR{Marketing}) & \currency{75,000.00} \\
مصاريف إدارية (\LR{Administrative}) & \currency{55,000.00} \\
إهلاك (\LR{Depreciation}) & \currency{40,000.00} \\
\midrule
\textbf{إجمالي المصاريف التشغيلية (\LR{Total Opex})} & \textbf{\currency{350,000.00}} \\
\midrule
\rowcolor{totalColor}
\textbf{\large الدخل التشغيلي (\LR{Operating Income})} & \textbf{\large\color{accentColor}\currency{330,000.00}} \\
\midrule
مصاريف تمويلية وضرائب (\LR{Financing \& Tax}) & \currency{60,000.00} \\
\midrule
\rowcolor{totalColor}
\textbf{\large صافي الدخل (\LR{Net Income})} & \textbf{\large\color{accentColor}\currency{270,000.00}} \\
\bottomrule
\end{tabularx}
\end{table}

\vspace{1.5em}

% ============================================================
% 2. BALANCE SHEET SUMMARY
% ============================================================
\section*{ملخص قائمة المركز المالي (\LR{Balance Sheet Summary})}
% TODO: replace sample amounts with actual figures

\begin{table}[H]
\centering
\renewcommand{\arraystretch}{1.4}
\begin{tabularx}{\textwidth}{L{10cm} R{5cm}}
\toprule
\textbf{البند (\LR{Item})} & \textbf{المبلغ (\LR{Amount})} \\
\midrule
\rowcolor{revenueColor}
\multicolumn{2}{L{15cm}}{\textbf{الأصول (\LR{Assets})}} \\
\midrule
الأصول المتداولة (\LR{Current Assets}) & \currency{420,000.00} \\
الأصول الثابتة (\LR{Fixed Assets}) & \currency{680,000.00} \\
أصول أخرى (\LR{Other Assets}) & \currency{100,000.00} \\
\midrule
\textbf{إجمالي الأصول (\LR{Total Assets})} & \textbf{\currency{1,200,000.00}} \\
\midrule
\rowcolor{expenseColor}
\multicolumn{2}{L{15cm}}{\textbf{الخصوم (\LR{Liabilities})}} \\
\midrule
الخصوم المتداولة (\LR{Current Liabilities}) & \currency{230,000.00} \\
الخصوم طويلة الأجل (\LR{Long-term Liabilities}) & \currency{370,000.00} \\
\midrule
\textbf{إجمالي الخصوم (\LR{Total Liabilities})} & \textbf{\currency{600,000.00}} \\
\midrule
\rowcolor{totalColor}
\multicolumn{2}{L{15cm}}{\textbf{حقوق الملكية (\LR{Equity})}} \\
\midrule
رأس المال (\LR{Capital}) & \currency{400,000.00} \\
الأرباح المرحّلة (\LR{Retained Earnings}) & \currency{200,000.00} \\
\midrule
\textbf{إجمالي حقوق الملكية (\LR{Total Equity})} & \textbf{\currency{600,000.00}} \\
\midrule
\rowcolor{totalColor}
\textbf{\large إجمالي الخصوم وحقوق الملكية (\LR{Total Liab. \& Equity})} & \textbf{\large\color{accentColor}\currency{1,200,000.00}} \\
\bottomrule
\end{tabularx}
\end{table}

\vspace{1.5em}

% ============================================================
% NOTES
% ============================================================
\section*{ملاحظات (\LR{Notes})}
\begin{itemize}[leftmargin=2em,itemsep=0.3em]
    \item جميع المبالغ بالريال السعودي (\LR{SAR}) ما لم يُذكر خلاف ذلك. % TODO: set currency
    \item تم إعداد هذه القائمة وفقاً لمعايير المحاسبة المعتمدة. % TODO: set accounting standard
    \item الأرقام تقريبية لأغراض العرض وقد تخضع للتعديل بعد المراجعة.
\end{itemize}

\vspace{1.5em}

\noindent
\begin{tabularx}{\textwidth}{@{} X X @{}}
\textbf{إعداد (\LR{Prepared By}):} & \textbf{اعتماد (\LR{Approved By}):} \\[2em]
\rule{6cm}{0.4pt} & \rule{6cm}{0.4pt} \\[0.5em]
{\small المسمى الوظيفي: \rule{3cm}{0.4pt}} & {\small المسمى الوظيفي: \rule{3cm}{0.4pt}} \\
\end{tabularx}

\end{document}
